\section{Counting on every subspace}\label{app:windows}

On $37$ of the $113$ intervals of Theorem~\ref{thm:profile-D} the corank bound comes from the counting
inequality applied to all constant subspaces at once. These intervals lie in $(3N,16.2N)$. No subspace is
constructed. The proof bounds the maximum over subspaces of the right side of Theorem~\ref{thm:everyU} by forced intersections of row spaces, and the
only finite data are upper bounds for ranks.

Let $V$ be a finite-dimensional vector space over $\F_p$. Let $c$ be a nonnegative integer. For $j$ in a finite
set $\mathcal J$ let $B_j\colon V\to W_j$ be a linear map and $w_j$ a nonnegative integer, and let $\lambda_{j,k}$ ($k<w_j$) be linear forms on
$\F_p^c$. Let $\mathcal C$ be the common kernel of the maps $B_j\circ(\lambda_{j,k}\otimes\mathrm{id})$ on $\F_p^c\otimes V$. A \emph{parametrization} is a linear map
$\iota$ from $\F_p^c\otimes V$ to the data $x$ of a decomposition modulo $p$ (Definition~\ref{def:slot}) with three properties on $\mathcal C$. The value
$\iota(x)$ is admissible, the map $\iota$ is injective on $\mathcal C$, and the row combination $c(\iota(x))$ satisfies the hypothesis of
Theorem~\ref{thm:count} with some polynomial of degree less than $r$. Nothing is required outside $\mathcal C$.

\begin{theorem}\label{thm:everyU}
Let $N\ge1$ and let $p$ be an odd prime with $H<p^2$. For a parametrization and every subspace $U\subset V$,
\[
 \nu_N(p)\ge c\dim U-\sum_{j\in\mathcal J}w_j\rank(B_j|_U).
\]
\end{theorem}

\begin{proof}
Since the coefficients $\lambda_{j,k}$ are scalars, $\lambda_{j,k}\otimes\mathrm{id}$ maps $\F_p^c\otimes U$ into $U$. So the restriction of $\iota$ to $\F_p^c\otimes U$ is
a parametrization with $V$ replaced by $U$ and $B_j$ by $B_j|_U$, and its common kernel is $\mathcal C\cap(\F_p^c\otimes U)$. Let $\Phi$ be the map on $\F_p^c\otimes U$ with
components $B_j\circ(\lambda_{j,k}\otimes\mathrm{id})$. These components map into $B_j(U)$, so $\rank\Phi\le\sum_jw_j\rank(B_j|_U)$. The
hypotheses of Theorem~\ref{thm:count} hold on $\ker\Phi$, with the injective map $x\mapsto c(\iota(x))$ on this kernel. The theorem gives
$\nu_N(p)\ge\dim\ker\Phi\ge c\dim U-\sum_jw_j\rank(B_j|_U)$.
\end{proof}

In the applications $c$ and the $w_j$ are affine functions of $N$ and $p$ on an interval of primes, while $V$ and
the $B_j$ are fixed. The space $\F_p^c$ is a space of polynomials of degree less than $c$ in the variable of a
decomposition modulo $p$. The forms $\lambda_{j,k}$ are Taylor coefficients at $X=-1,0,1$ and coefficients at $\infty$. The map $\iota$ places the data block by block along the rows. Its three properties are proved once for the whole family of parametrizations. On $\mathcal C$ the moment columns take the
values of a polynomial of degree less than $r$. Since the intervals lie in $(3N,16.2N)$, where $p>q$, we have $a=0$ in
Theorem~\ref{thm:count}.

\emph{Lower bounds for the maximum.} For a subspace $S$ of the dual $V^*$ let $U=S^\perp$. Let $R_j\subset V^*$ be the row space of $B_j$, and let $r_j\ge\dim R_j$ be real
numbers, in practice upper bounds for the ranks. Since $\rank(B_j|_U)=\dim R_j-\dim(R_j\cap S)$, the right side of Theorem~\ref{thm:everyU} is
\[
 E(S)=c(\dim V-\dim S)-\sum_jw_j\bigl(\dim R_j-\dim(R_j\cap S)\bigr).
\]
Let $e_0=c\dim V-\sum_jw_jr_j$ and $\delta_j=r_j-\dim R_j\ge0$. Then
\begin{equation}\label{eq:eulerid}
 E(S)=e_0+\sum_jw_j\delta_j+\sum_jw_j\dim(R_j\cap S)-c\dim S .
\end{equation}
A map may also be omitted from the sum if its weight is subtracted from $c$, because $\dim R_j-\dim(R_j\cap S)\le\dim V-\dim S$. The following three lemmas give a subspace
with $E(S)>e_0$, from relations among the row spaces. All sets of indices are subsets of $\mathcal J$, and all $w_j$ are nonnegative.

\begin{lemma}[pair]\label{lem:pairU}
Let $a\ne b$ lie in a set $I$ such that $R_a\le R_j$ or $R_b\le R_j$ for all $j\in I$. Let
$s\ge\dim(R_a+R_b)$ and $0\le h\le\min\bigl(w_a,w_b,\sum_{j\in I}w_j-c\bigr)$. Then some $S$ has
$E(S)\ge e_0+\max(0,r_a+r_b-s)\,h$.
\end{lemma}

\begin{proof}
Take $S=R_a\cap R_b$ and let $k=\dim S$. Then $S\le R_j$ for $j\in I$, so \eqref{eq:eulerid} gives
\[
 E(S)-e_0\ge w_a\delta_a+w_b\delta_b+\Bigl(\sum_{j\in I}w_j-c\Bigr)k\ge h(\delta_a+\delta_b+k).
\]
Since $\dim(R_a+R_b)+k=\dim R_a+\dim R_b$, we have $r_a+r_b-s\le\delta_a+\delta_b+k$. Also $\delta_a+\delta_b+k\ge0$.
\end{proof}

\begin{lemma}[several points]\label{lem:pointsU}
Let $c\ge0$. Let $I_1,\ldots,I_l$ be disjoint sets with $a_i\in I_i$, such that $R_{a_i}\le R_j$ for $j\in I_i$. Put $W_i=\sum_{j\in I_i}w_j$ and suppose that $W_i\le c$ for each $i$.
Let $s\ge\dim\sum_iR_{a_i}$. Let $h\ge0$ with $h\le w_{a_i}$ for each $i$ and $h\le\sum_iW_i-(l-1)c$. Then some $S$ has
$E(S)\ge e_0+\max(0,\sum_ir_{a_i}-s)\,h$.
\end{lemma}

\begin{proof}
Let $\mathsf T$ be the kernel of the sum map $\bigoplus_iR_{a_i}\to V^*$, and let $k=\dim\mathsf T$. Let $K_i\le R_{a_i}$ be the projection of $\mathsf T$ to the summand $i$, let $\kappa_i=\dim K_i$, and let
$S=\sum_iK_i$. Then $\kappa_i\le k$. The rank formula for the sum map gives $\sum_i\dim R_{a_i}\le s+k$. The kernel of the sum map on $\bigoplus_iK_i$ is again $\mathsf T$, so
$\dim S=\sum_i\kappa_i-k$.

For $j\in I_i$ we have $K_i\le R_j\cap S$, so $\dim(R_j\cap S)\ge\kappa_i$. Since the $I_i$ are disjoint, the $a_i$ are distinct, and \eqref{eq:eulerid} gives
\[
 E(S)-e_0\ge\sum_iw_{a_i}\delta_{a_i}+\sum_iW_i\kappa_i-c\sum_i\kappa_i+ck\ge h\sum_i\delta_{a_i}+\sum_i(W_i-c)\kappa_i+ck.
\]
Since $W_i-c\le0$ and $\kappa_i\le k$, the last two terms are at least $k(\sum_iW_i-lc+c)\ge hk$. So $E(S)-e_0\ge h(\sum_i\delta_{a_i}+k)$, and
$\sum_ir_{a_i}-s\le\sum_i\delta_{a_i}+k$.
\end{proof}

\begin{lemma}[kernel]\label{lem:kernelU}
Let $R\subset V^*$ be a subspace with $\dim R\le s$, and let $Z$ be a set with $R_j\le R$ for $j\in Z$. For every integer $d\ge0$
with $d+s\le\dim V$ some $S$ has
$E(S)\ge cd-\sum_{j\notin Z}w_j\min(r_j,d)$.
\end{lemma}

\begin{proof}
The space $R^\perp\subset V$ has dimension at least $d$. Choose $U\le R^\perp$ of dimension $d$ and let $S=U^\perp$. Then $B_j|_U=0$ for $j\in Z$, and $\rank(B_j|_U)\le\min(r_j,d)$ for
the other $j$.
\end{proof}

\emph{The rank data.} The space $V$ is a space of profiles of the coefficient vectors of a block, with coordinates
$T_k$. The conditions at the moment columns are linear combinations of a small number of basic rows
\[
 A[k]\mapsto T_k+T_{k+1},\qquad C[k]\mapsto T_{k+1}-T_k,
\]
together with the row of the centre coordinate. So the moment values themselves are not needed. Every row factors
through the map $G$ whose components are these basic rows. This gives $r_j\le\min(\rank F_j+\#\text{basic rows},\ \rank(G,F_j))$, where $F_j$ is the fixed part
of the map. Each rank bound is certified by an integer left inverse, checked by the Lean kernel. Integer left
inverses are valid at every prime, so no prime enters this step.

The $37$ intervals are treated as follows: $13$ by Lemma~\ref{lem:kernelU}, $12$ by Lemma~\ref{lem:pointsU}, $11$ by
Lemma~\ref{lem:pairU}, and one by $E(V)$ itself. On each interval the bound is an affine function
$\nu_N(p)\ge c_0N+c_1p-2$ with rational $c_0$ and $c_1$. The rounding of the integer orders is absorbed by
the shortening $\tau$ of the interval and the lowering of the bound by $10^{-6}$ in Section~\ref{ssec:kappa}.

\section{The bounds for the energy method}\label{app:field}

This appendix gives the trial sequences of Section~\ref{sec:energy}, the certified bounds that they satisfy, and
the bounds for the weights of Section~\ref{sec:real}.

\emph{Trial sequences.} For the constrained case we use a trial sequence with $40$ Chebyshev coefficients and a
finite number of geometric terms, with the multiplier $\lambda_1=1.68493357$. For the case with doubled
interaction we use trial sequences of the form \eqref{eq:trial}, with $\Kcoef$ Chebyshev coefficients, which are
dyadic rationals, and $\Ngeo$ geometric terms. Of these, $\Ncomplex$ have complex $z_j$ and complex $r_j$. We
place the complex charges at the $41$ equally spaced points $x_0=-0.16+0.008j$ $(0\le j\le40)$, all with
$\epsilon_0=10^{-2}$. The real terms sit at points $\pm z$ with $z$ increasing from $0.8$ to $0.999997$. The
numbers $z_j$ are rationals with at most twelve decimals. The norms of the geometric terms have closed forms:
two terms $(z_j,r_j)$ and $(z_l,r_l)$ contribute
$-\frac12\operatorname{Re}\bigl(r_jr_l\log(1-z_jz_l)+r_j\bar r_l\log(1-z_j\bar z_l)\bigr)$ to $\|u\|^2$.

\emph{The suprema of the fields.} We must bound two functions of one variable from above on their domains: the
field $X$ on $[-1,1]$ and the field $Y$ on $[0,1]$. We locate the real
critical points. The derivative is a rational function with a known denominator, and an exact identity gives its
numerator. The Descartes rule, after a M\"obius transformation of the interval, bounds the number of real roots
of the numerator. Sign evaluations at $81$-bit rational points then isolate this number of roots. One global Lipschitz
bound on each isolating interval and the values at the end points give the supremum.

\emph{The bounds.} For the base trial sequences this gives
\[
 \sup X\le\SupX,\quad \sup Y\le\SupY .
\]
So $M_X+M_Y=\FieldCap$, and the norms satisfy $2\|u\|^2+\frac12\|v\|^2\le0.7736206$. The first coefficient satisfies
$u_1\ge9/50$, which is used near $x=0$. For the constrained case, with the weights
$(\alpha_A,\beta_h,\gamma_H-\frac12,\lambda_1,2g_0/n_0)=(11/29,3/29,49/58,\lambda_1,5/29)$, the bounds
$\sup X_1\le-1.15854$ and $\sup Y\le-1.32981$ are obtained by isolating the zeros of the numerators of the
derivatives, which are polynomials with integer coefficients. The norm bound is $n_1=0.67761$.

\emph{The increments.} The increment of the case with doubled interaction has $40$ Chebyshev coefficients, and
that of the constrained case has $24$. Both are dyadic. Their fields are bounded on the whole domains by the
same methods. For the case with doubled interaction the norm bound at $\tau$ is $0.7736206-0.117647903\,\tau+0.005215061\,\tau^2$,
the field $X$ gains at most $0.01484272\,\tau$, and the column field $Y$ with $\gamma g$ is bounded by a piecewise
linear function of $\gamma$ with knots at $\alpha_1,\frac15,\frac3{10},\frac25,\frac35,\frac45,1$. The fields $F_1$ and $F_2$
contain the node term $\hat\delta/58$. They are bounded by covers by cells, with exact rational and integer arithmetic on each cell
and the maximum of $\hat\delta$ on each cell. For $F_1$ on $|x|\ge401/2500$ the cover has $150$ cells. It uses a Taylor expansion of the
polynomial part and an expansion to order $5$ of the logarithmic terms. On the positive part of $E_W$ below $401/2500$ six
further cells suffice. For $F_2$ the covers of $[98749/10^6,\,9995/10^4]$ have $17$ and $14$ cells at the two ends of the move. This gives
\[
 \sup F_1\le0.007,\qquad \sup F_2\le0.007-\log2
\]
at $\tau=0$, and $-0.0268$ and $0.0233-\log2$ at $\tau=1$. For the constrained case the norm bound at $\tau$ is
$0.67761-0.082834152\,\tau+0.004096781\,\tau^2$, and the bound $\mathcal M_1$ of Section~\ref{sec:energy} is the chord of
$0.35072$ and $0.35295$. The two numbers are maxima over the $865$ segments of the chord table at the two end values of
$\lambda_1$. With $a_0=41/58$, $a_1=76/29$ and $U_0=\Uzero$ these bounds give the inequalities of \S\ref{ssec:cases}.
Each check $V(\alpha)\le\sigma\alpha$ is a polynomial inequality on a closed interval and is proved by a cover in exact
rational arithmetic.

\emph{The bound for the weights.} Inequalities \eqref{eq:circlecap} and \eqref{eq:segcap} concern the explicit function $L$ of
\eqref{eq:Ldef}. On the circle $|z|=\frac34$ write $z=\frac34(s+i\sqrt{1-s^2})$ with $s=\cos\theta$. Then $|z-t|^2=\frac9{16}+t^2-\frac32ts$
is affine in $s$ for every real $t$. So $\log|1-z|$ is a concave function of $s$. Each $G_c(z)$ is $\operatorname{arccosh}$ of
$(|z-a_c|+|z-b_c|)/w_c$, a sum of square roots of affine functions of $s$, and therefore concave in $s$ as well. Hence the left
side of \eqref{eq:circlecap} is a concave function of $s\in[-1,1]$. Its maximum is $3.25053\ldots$ near $s=-0.1741$. A concave function
is bounded by its secants outside the points and by the chord between them, so three points at distance $\frac1{2000}$ around the
maximum suffice. At these points the value is enclosed to within $10^{-17}$ by exact rational arithmetic with integer square roots,
and the bound $D_\Gamma=\Dcirc$ holds with a margin $6\cdot10^{-5}$. On the real segments we use, for real $x$, the closed form
$e^{G_c(x)}w_c/4\le\frac12\max(w_c,|2x-a_c-b_c|)$ together with $6\log|1-x|\le6\log\frac14$ for $x\ge\frac34$ and
$6\log|1-x|\le6\log2$ for $x\le-\frac34$. This gives \eqref{eq:segcap} with values below $-9$, far below $D_\Gamma$.
The knot values $v_i$ of the chord table lie between $-1.97$ and $17.99$.

\section{The reference measures of Section~\ref{sec:signed}}\label{app:trials}

The bound of \S\ref{ssec:budget} uses reference measures $q_+$ and $q_-$, a row reference $R$ and a column reference $C$.
This appendix describes them and the certified bounds \eqref{eq:caps}.

\emph{The measures.} On the real line of the variable $Y=x+1/x$ with $|Y|\ge2$ we take measures $G_+$ and $G_-$ that
are piecewise uniform on finitely many cells in $Y$, with masses $m_+$ and $m_-$ and total mass $1$. They carry the positive and the
negative free nodes, and $R=G_++G_-$ is the row reference. The column reference $C$ has mass $1$. The measures are completed by
auxiliary measures on the $16$ inner circles. Each auxiliary measure has a cubed Poisson density: its Fourier coefficients
are $-q^k(1+Ak+Bk^2)/k$ with $q=e^{-\varepsilon/2}$ and $\varepsilon=27/250$, and the constants $A$ and $B$ are rational functions of $r=q^2$.
The auxiliary measures $\nu_+$ and $\nu_-$ sit at the phases $\pi j/93$, with $616$ and $935$ entries respectively, and have total masses $1-m_+$ and $\frac{53}{58}-m_-$. Let $X$ be the inverse of $x\mapsto x+1/x$ with $|X|\le1$. Then
$q_+=X_\#G_++\nu_+$ is the measure in the plane of $z$. In the plane of $u$ we take $q_-=J_\#\tilde q_-$ with $\tilde q_-=X_\#G_-+\nu_-$.

\emph{The fields.} Let $L(\mu)=\int\log|X|\,d\mu$. For $\sigma=\pm$ the free row field is
\begin{multline*}
 F_\sigma(x)=L(R)+L(G_\sigma)-\kappa_\sigma+U^{G_\sigma+R-2_\#C}(x+1/x)+\bigl(m_\sigma+\tfrac{11}{29}\bigr)\log|x|\\
 -\tfrac{39}{29}\log(1+x^2)+\tfrac3{29}\log(1-x)+[\sigma=+]\bigl(\mathsf P(x)+\mathsf E\bigr)+A_\sigma(x),
\end{multline*}
where $A_+=U^{\nu_+}$ and $A_-(x)=\int(\log|x-w|+\log|1-xw|)\,d\nu_-(w)$. The polynomial $\mathsf P(x)=\sum_{k\le900}m_kx^k/k$ with $m_k=\int X^k\,dG_+$ is
the truncated reflection potential of the positive measure, and $\mathsf E\ge0$ bounds its tail.
The column field is
\[
 \mathsf C(s)=-L(R)+2U^C(s)-U^R(2s)+\tfrac9{58}\log s+\tfrac3{58}\log(1-s).
\]
A prescribed row outside $S$ has the field $\mathsf H_\sigma$ obtained from $F_\sigma$ by three changes. The references $R$ and $C$ are removed, the
coefficients $\frac{11}{29}$ and $\frac3{29}$ are replaced by $0$, and $\frac{39}{29}$ is replaced by $\frac{53}{58}$. Its domain is
$x\le-19501/10^5$ and $x\ge98749/10^6$, which contains $E_x$. A prescribed row in $S$ has the field
\[
 \mathsf H_S(y)=\tfrac{48}{58}\log y-\tfrac{53}{58}\log(1+y^2)-\kappa_+-\mathsf P(y)+\int\log|1-yw|\,d\nu_+(w)
\]
for $0<y<1$.
The angular fields are the same potentials on the $16$ circles, together with the terms $(53\sigma_j-\tau_j/2)/58$.

\emph{The constant $\mathsf I$.} Let $E(\mu,\nu)=\iint\log|Y-Y'|\,d\mu\,d\nu$. Evaluating $\mathsf Q$ on the references gives $\mathsf I_0=-\frac12E(R,R)-E(C,C)+E(R,2_\#C)$. Since $R=G_++G_-$, six pair
energies of lists of cells occur: those of $G_+$, $G_-$ and $C$ with themselves, $E(G_+,G_-)$, $E(G_+,2_\#C)$ and $E(G_-,2_\#C)$. The energies $I(q_\pm)$ split into cell, cross and auxiliary parts. Their cell parts contain the same
pair energies. For piecewise uniform measures a pair energy is a bilinear form in the jumps of the densities, with the kernel $t^2\log|t|$, up to constants. We evaluate it on integer end points in fixed point arithmetic. Well separated blocks of end points are
treated by a series in the offsets from their centres, using the moments of the blocks, and the other pairs by a recursion on the blocks. This encloses each pair energy to within $10^{-7}$. The remaining terms of $I(q_\pm)$ are the logarithmic moments of the cells,
the integrals of the cells against the rings, which are sums of primitives over the end points, and the quadratic forms of the auxiliary measures, which have closed forms on the rings. All of these are enclosed to within $10^{-9}$. The reflection term $-\iint\log(1-xx')$ of $I(q_+)$
enters only through a lower bound from its truncated series. Finally $\int_0^{53}\tau^*\ge1425.3339$. Together these give $\mathsf I_0+\mathsf I_{\mathrm{gas}}\le2.366997$.

\emph{The column field.} The three bounds $c^{\mathsf C}_i$ of \eqref{eq:caps} on $(0,1)$ use a chain of $386$ units $s=(M+\delta z)/10^{12}$ with $|z|\le1$, together with two end ranges. On a unit the field is at most an integer polynomial in $z$ plus an explicit
error. The far cells of the references enter through exact integer moments of contiguous clusters, $1152$ clusters in all, and the Taylor series of the logarithmic kernel. The near cells of $C$ enter through the closed form of their potential. The two logarithms enter by
series or, on the end units, by monotonicity. The constant is $L(R)\ge-0.6905432$, from the primitive of $\log|X|$ in $Y$.

\emph{The angular bounds.} The bounds $\kappa_\pm$ are proved for both signs on each of the $16$ radii, by $32$ certificates. The auxiliary part is cut at degree $384$, and its tail is bounded. The gas is split into bins. On a bin the kernel has the form
$\mathsf Q_u(x)=\mathsf a-2\mathsf bux+\mathsf cx^2$ with $u=\cos\theta$. Jensen's inequality and an odd truncation of the logarithmic series bound $\int\frac12\log\mathsf Q_u$ by a polynomial in $u$. The negative sign has a direct and a reflected kernel. The sum over the bins
is converted exactly to a Chebyshev series. Its bound by $\kappa_\pm$ on $[-1,1]$ is certified by an identity of the same form as the one for the free rows below.

\emph{The prescribed rows.} The auxiliary series of $\mathsf H_\pm$ are cut at degree $384$. The gas potential is expanded in Chebyshev polynomials on $[0,1]$. The first $64$ terms are exact, and the tail is bounded by Bessel's inequality from a bound for the $L^2$ energy of the cell densities.
The reflection series is dropped by sign for the positive gas and kept through exact moments for the negative gas. The radial field is a sum of $16$ blocks, one for each radius. Its polynomial part goes exactly into the Chebyshev variable, and the part of degree at most $16$ is bounded by
a chord cover with the radii as nodes. This gives a polynomial majorant, and a sum of squares identity gives $\mathsf H_\pm\le-0.9783$. On $[\frac1{10},1]$ the field $\mathsf H_S$ is treated in the same way. There $\mathsf P$ is replaced by its partial sum of degree $32$, since the moments are nonnegative, the auxiliary series is cut at degree $32$ with an absolute tail, and
$\log y$ and $\log(1+y^2)$ are expanded explicitly. This gives $\mathsf H_S\le-0.997$. On $(0,\frac1{10}]$ one has directly $\mathsf H_S(y)\le\frac{24}{29}\log y+\log(1+y)-\kappa_++$(a fixed tail), which is smaller. So $\mathsf M=-0.976$ is admissible.

\emph{The free rows.} The free row fields are bounded on each sign as follows. Put $t=|x|$ and $z=2t-1$. The exact field is at most its truncation to degree $384$ in the auxiliary terms, up to a tail
below $10^{-8}$. We bound the truncated field by one polynomial in $z$. The direct potential of the own measure, $\int\log|t-X(y)|$, is replaced by its Chebyshev series,
with a tail bounded through the variation of the density. Thin pieces of nonnegative density are dropped, since $\log|t-v|\le0$ there. All other potentials
are bounded from above by finite sums of terms $w\log|1-r(\mathsf q)t|$ with rational $\mathsf q$: tangent majorants for the reflection and opposite pieces, chord majorants for the
column panels, and the same form for $\log t$ and $\log(1-t)$. Each such term has an explicit Chebyshev expansion with a geometric tail. The moments of these terms and the
direct moments are computed by exact integer arithmetic with Kronecker packing. The radial field and the term $-\frac{39}{29}\log(1+t^2)$ are bounded by a chord cover and a
Chebyshev polynomial. This gives a polynomial $\mathsf p(z)$ with $F_\sigma\le\mathsf p$ plus explicit small errors. Finally $\mathsf p(z)\le\mathsf c$ on $[-1,1]$ is proved by an exact identity
$\mathsf d(\mathsf c-\mathsf p)=\mathsf r+\mathsf q_0^2$ between Chebyshev series, with $\mathsf q_0$ of degree $16383$ and with $\mathsf r$ small in the $\ell^1$ norm of its coefficients. This gives
the bounds $c^{\mathsf R}_0$ and $c^{\mathsf R}_1$ of \eqref{eq:caps} at the two ends of the move, for both signs.
